\chapter{Project management}
\label{chap:project-management}

This chapter reflects on the management strategies employed throughout the project lifecycle. Our team adopted an Agile-inspired approach, utilizing a Kanban workflow to manage the transition from a client-led project to an academic one.

\section{Team Structure and Roles}
To ensure accountability and specialization, each team member was assigned a specific leadership role alongside their development duties \cite{uploaded:Key points for report .docx}:
\begin{itemize}
    \item \textbf{Hamed Latif}: Team Lead, Client Liaison, and Final Submission Manager
    \item \textbf{Kausar Khan}: Trello Board Manager and Note-taker
    \item \textbf{Eric Almeida}: User Story and Task Manager
    \item \textbf{Aishah Ahmed}: LaTeX Report Manager
    \item \textbf{Cate Tay (Cate) Chew}: GitHub and Version Control Manager
    \item \textbf{Zerun (Harry) Chen}: Deployment Manager
    \item \textbf{Dian Mahmood}: Code Quality Manager
    \item \textbf{Indixora Jozal}: People’s Engagement Manager
    \item \textbf{Binuli Nevanma Markawitage}: Meeting Setup Manager
\end{itemize}

\section{Project Evolution}
The project began as a client-facing application for \textit{Makcorp}, focused on analyzing ASX mining stocks. However, due to significant delays and poor communication channels during the first two weeks, the team collectively decided to pivot to an academic project: \textbf{StudyFlax}, a time management system for students. 

To mitigate the risk of reduced development time following this switch, the team increased its target goal to 18 user stories per week. This allowed us to regain momentum and ensure all core features of StudyFlax were delivered on schedule.

\section{Communication and Collaboration}
Our communication strategy relied on a combination of synchronous and asynchronous tools:
\begin{itemize}
    \item \textbf{Meetings}: We held two formal meetings per week (Mondays at 3:00 PM and Fridays at 2:00 PM) to discuss progress and blockers.
    \item \textbf{Instant Messaging}: WhatsApp was used for daily coordination, with a protocol requiring members to reply within the same day.
    \item \textbf{Task Tracking}: A Trello Kanban board was used to visualize the backlog, in-progress tasks, and completed features.
\end{itemize}

\section{Software Quality and Standards}
To maintain a high standard of software engineering, the team enforced strict coding conventions:
\begin{itemize}
    \item \textbf{Code Reviews}: All pull requests required approval from at least three team members before merging into the main branch.
    \item \textbf{Function Length}: Methods were limited to a maximum of 15 lines to ensure readability and low coupling.
    \item \textbf{Testing}: We aimed for 100\% code coverage, with testing methods capped at 25 lines .
\end{itemize}

\section{Reflection and Lessons Learned}
The most significant lesson learned was the importance of early risk identification. While the initial communication issues with Makcorp were a setback, communication within the team and a clear division of responsibilities ensured the team remained cohesive during the transition. The use of a decentralized leadership structure allowed us to move faster than a traditional hierarchy would have permitted. For future projects, we would establish a more robust contingency plan for client-side delays earlier in the process.